% =============================================================================
%  Appendix — core code
%  顺序：先附件说明表，再贴代码。附表说明"哪个文件对应哪个问题、产出什么"，
%  审查者找不到映射，就不会去读代码。
%  篇幅：附录压在正文的 1/3 以内；放不下就把完整代码单独提交，正文只留核心循环。
% =============================================================================

\subsection*{A.1 File list}

\begin{table}[htbp]
  \centering
  \caption{Submitted code files.}
  \label{tab:codefiles}
  \footnotesize
  \begin{tabular}{@{}p{0.17\textwidth}p{0.15\textwidth}p{0.40\textwidth}p{0.16\textwidth}@{}}
    \toprule
    File & Problem & Purpose & Output \\
    \midrule
    \texttt{utils.py}    & all       & parameters, index evaluation, plotting & -- \\
    \texttt{problem1.py} & Problem 1 & [what it does]                         & \texttt{[file]} \\
    \texttt{problem2.py} & Problem 2 & [what it does]                         & \texttt{[file]} \\
    \bottomrule
  \end{tabular}
\end{table}

\subsection*{A.2 Core code}

\begin{lstlisting}[language=Python]
# problem1.py -- forward simulation
# run: python problem1.py   -> writes [output file]
import numpy as np
from utils import params, vibration_index

P = params()
# [core loop: one step per sample, accumulate the index]
\end{lstlisting}

% 【Appendix writing prompt】
% 每个 listing 第一行写死 "run: ... -> output"，读者不用猜怎么跑。
% 一题一份可运行脚本；草稿版、重复 import 的副本不要贴。
